\documentclass[11pt]{article}
\usepackage{explorer}
\newcommand{\ExplorerLanguage}{FR}
\newcommand{\ExplorerNumber}{09}
\newcommand{\ExplorerKind}{Problème de synthèse}
\begin{document}
\ExplorerTitle{Mathématiques expertes · géométrie}{Rotations et nombres complexes}
\ExplorerSection{Chercher, justifier, conclure}
On définit \(z_{0}=2\) et \(z_{n+1}=i z_{n}+2\). Trouver le point fixe a, puis une formule explicite de \(z_{n}\). Déterminer la période minimale de la suite et la somme \(z_{0}+z_{1}+z_{2}+z_{3}\). Interpréter géométriquement les quatre points.\par
\vspace{5pt}
\end{document}
